% ==========================================================================
% GENERATED by scripts/split_paper.py. Do not edit this file.
%
% Source: imported_papers/Koopman-JEPA-TMLR-rewrite/main.tex
% Part 5: Results
%
% Edits belong upstream, in the published paper. To change what the
% chapter includes, change imports/bounded.toml and re-run
% the splitter.
% ==========================================================================
\section{Results}
\label{sec:exp}

\subsection{Measured the Standard Way, the Constraint Wins}
\label{sec:results}
\input{tables/table_main_controls.tex}
% [orig 5.1 + new] DONE. The first paragraph keeps the approved paper's
% reading of the separation, in its own terms, on the seeded numbers. The
% second is new and is the point of the section.
Measured the way this literature measures, the constraint works. Table~\ref{tab:combined} reports probe state MSE at $T{=}10$ and $T{=}20$ on clean backgrounds at full data. At $T{=}20$ on pendulum the four dynamics modules span more than five orders of magnitude, from $3.68$ for Koopman~(Stable) to $8.6\times10^{5}$ for the unconstrained operator, with the MLP-JEPA baseline at $6.4\times10^{3}$. Cartpole separates the same way and a constrained operator again leads, Koopman~(Soft) at $13.5$ against $1.9\times10^{3}$ for the baseline. The spectra are consistent with the ordering: the unconstrained operator's eigenvalues drift outside the unit circle ($\rho(\K) = 1.43$), the soft variant settles just outside it ($1.04$), and the Cayley variant sits exactly on it by construction (Figure~\ref{fig:eigenvalues}). Linearity by itself is not what separates them, since the unconstrained linear operator diverges alongside the nonlinear baseline.

The control rows change the reading. One of the two trivial predictors is the most accurate entry in every block of the table. On pendulum at $T{=}20$ the best learned model reads $3.68$ against $1.66$ for copying the current latent forward. On cartpole it reads $13.5$ against $0.98$ for the dataset mean. Counting every environment, dynamics module and horizon at $T \ge 10$, a trivial predictor wins $24$ of $24$ cells, and $168$ of $168$ once all four backgrounds and all four data fractions are included (Appendix~\ref{app:full_results}). The learned models win only at the shortest horizon, and only on cartpole, where three of the four beat the floor by between $4$ and $15$ percent.

Section~\ref{sec:bound_limits} anticipated this. A predictor that cannot amplify its input has a flat error curve whether or not it has learned the dynamics, and both controls are of that kind. What Table~\ref{tab:combined} establishes is therefore not that the constraint helps, but that the probe ranks the constrained operators first while ranking two predictors that learned nothing ahead of all of them.

\subsection{The Advantage Evaporates under a Scale-Invariant Readout}
\label{sec:inversion}
\input{tables/table_retrieval.tex}
% [draft 5.2, rewritten on the seeded numbers] DONE. Two claims from the draft
% are gone because the data does not support them. The draft said the ranking
% inverts; it does not, the order is identical under both readouts on both
% environments and what collapses is the spread. The draft also said the two
% readouts agree on persistence to within one percent; the measured gap is four
% percent on pendulum and nine on cartpole, so the sentence says "a few
% percent". The decoder sentence is dropped: visual MSE was only ever measured
% on the superseded pass and is not in the ledger.
Table~\ref{tab:retrieval} scores the same rollouts twice, once with the probe and once with the scale-invariant readout of Section~\ref{sec:measurements}. Under the probe at $T{=}20$ on pendulum the four methods span a factor of $2.3\times10^{5}$, from $3.68$ for Koopman~(Stable) to $8.6\times10^{5}$ for the unconstrained operator. Under retrieval the same rollouts span a factor of $1.38$, from $4.05$ to $5.61$. On cartpole a spread of $139$ becomes $1.50$. Nothing about the rollouts changed between the two columns, only what is read off them.

The ordering does not change. Under both readouts, on both environments, the four methods rank in the same order, so this is a collapse rather than a reshuffle: a separation of five orders of magnitude becomes a separation of forty percent. Nor does the scale-invariant readout reveal a hidden winner. No learned model is meaningfully better than persistence under it, the closest being the soft variant on cartpole at $1.12$ against $1.13$.

The ratio column localizes the disagreement. It is near one for persistence, which by construction neither grows nor shrinks the latent, so the two readouts agree on it to within a few percent. It is $0.91$ for Koopman~(Stable), whose operator preserves norm by construction, rises to $4.3$ for the soft variant, whose spectral radius sits slightly above one, and reaches $1.3\times10^{3}$ for the baseline and $1.5\times10^{5}$ for the unconstrained operator. That is the ordering of how far each operator lets the latent norm travel, which is the quantity Section~\ref{sec:variants} noted the training objective does not constrain.

That last step is an inference from two readouts, and it can be checked directly on the operators. Rolling encoded held-out frames forward and recording the median of $\|\K^k\z\| / \|\z\|$ gives, on pendulum at $T{=}20$, exactly $1.000$ for Koopman~(Stable), $4.2$ for the soft variant, $6.9\times10^{2}$ for the baseline and $6.0\times10^{3}$ for the unconstrained operator, and the ordering matches the ratio column (Appendix~\ref{app:drift}). Cartpole is a weaker test of the same claim. Every operator there stays within a factor of ten at $T{=}20$, and at that scale the growth ordering no longer reproduces the ratio ordering, as one should expect once the drift is small enough for other differences to matter.


\subsection{The Metric Measures the Constraint}
\label{sec:spectrum}
\begin{figure}[t]
    \centering
    \includegraphics[width=\linewidth]{paper_figures/fig_rho_sweep.pdf}
    \caption{Probe error is governed by the spectral radius, and the response
    is monotone in all six curves. Each trained operator is rescaled to a target
    radius $\rho$ and the rollout is run again with the encoder, the probe and the
    action matrix untouched, so the only quantity that changes along a curve is
    the spectrum. \textbf{Left:} pendulum-swingup. \textbf{Right:}
    cartpole-swingup. Probe state MSE at $T{=}20$, mean over 10 seeds, log
    scale. The dotted line marks $\rho = 1$, the value the Cayley
    parameterization enforces. The lowest error on every curve is at
    $\rho = 0.9$. Values are comparable within a curve and not against
    Table~\ref{tab:combined}, which evaluates more episodes.}
    \label{fig:rho_sweep}
\end{figure}
% [new 5.3] DONE. The rho=1.0 point is not the Table 1 value for the same
% operator, because the sweep evaluates 20 episodes and the Table 1 rows
% evaluate 50. Every comparison here therefore stays inside one curve, and the
% caption says so.
Rescaling a trained operator to a target spectral radius and re-running the rollout moves the probe error monotonically, in all six curves of Figure~\ref{fig:rho_sweep}. On pendulum the Cayley operator rescaled to $\rho = 0.9$ reads $2.75$ at $T{=}20$, and the same operator at $\rho = 1.2$ reads $2.3\times10^{3}$, a factor of $850$ from a change that touches nothing but the spectrum. The encoder, the probe, the action matrix and the learned features are identical along each curve.

This is the intervention that the four-way comparison cannot perform. Comparing architectures varies the spectrum together with everything else that differs between them, whereas rescaling varies the spectrum alone, and the response is what Equation~\ref{eq:linear_bound} implies. Error growth tracks $\rho^k$. The theory is not what fails here. Every curve rises where the bound says it should.

The consequence is the point of this paper. The probe reads a latent by mapping it to physical coordinates, and it was fitted at the scale of encoded observations, so it charges a rollout for letting its norm run away. The bound limits exactly how far a rollout's norm can run away. Probe error and the quantity the bound constrains are therefore two measurements of the same thing, and an operator that satisfies the bound is guaranteed a flat probe curve whether or not it has learned the dynamics. The agreement between Table~\ref{tab:combined} and Equation~\ref{eq:linear_bound} that the study was designed to test is not a confirmation. It is one quantity measured twice. A bounded error curve under a norm-sensitive readout is a property of the constraint, not evidence of a learned model.

Two further readings follow from the figure. Contraction beats norm preservation, since the lowest error on all six curves is at $\rho = 0.9$ rather than at the $\rho = 1$ the Cayley parameterization enforces, so the architecture fixes the spectrum at a value the data does not prefer. And at a common $\rho$ the three operators differ far less than any one of them differs across $\rho$, which is another way of saying that the four-way comparison measured mostly the spectral radius.

\subsection{A Scale-Aware Loss Removes the Gap}
\label{sec:loss}
\input{tables/table_l2_ablation.tex}
% [new 5.4] DONE. Both arms are trained for this ablation, so the cosine
% column is not the Table 1 column: it is fifteen fresh runs on pendulum at
% clean backgrounds and full data. Saying so matters, because the alternative
% reading is that the comparison mixes the ablation against the main study.
Section~\ref{sec:inversion} attributes the gap between the two readouts to a training objective that does not constrain latent norm. If that is right, an objective that does constrain it should close the gap. Table~\ref{tab:l2_ablation} reports the test. Two arms of five seeds each are trained on pendulum with clean backgrounds and full data, differing only in the prediction loss, and both are evaluated on the same held-out episodes, so the dataset-mean control reads $2.51$ at $T{=}20$ for both.

The gap closes. Under the cosine objective the probe reads $2.6\times10^{6}$ times worse than retrieval for the MLP-JEPA baseline at $T{=}20$, and $5.3\times10^{4}$ times worse for the unconstrained operator. Replacing the cosine distance with an $\ell_2$ prediction loss brings those ratios to $2.33$ and $0.55$. The Cayley operator is the control here rather than the test: its ratio is already near one because the operator preserves norm by construction, and it stays near one under both objectives, $0.895$ against $1.40$. Four orders of magnitude of disagreement between two readouts are a property of the loss, and they can be removed by changing it.

Removing them does not produce a better model. Under the $\ell_2$ objective a trivial predictor is still more accurate than every learned model in all six cells, by $9.89$ against $1.65$ for the baseline at $T{=}20$ and $3.21$ against $1.71$ for the unconstrained operator. The recommendation this supports, which Section~\ref{sec:checklist} states, is about what the metric measures rather than about how well the models do.

\subsection{The Benchmark Degenerates, and the Control Claim Does Not Transfer}
\label{sec:degeneracy}
\begin{figure}[t]
    \centering
    \includegraphics[width=0.48\linewidth]{paper_figures/fig_degeneracy.pdf}
    \caption{The pendulum benchmark stops discriminating at long horizon and
    cartpole does not. Error of the dataset-mean predictor against rollout
    horizon, log scale. A constant predictor cannot improve as the horizon
    grows unless the states it is asked to predict are themselves becoming more
    predictable, so the falling pendulum curve is evidence that they are. The
    predictor uses no model, so each curve is a property of the evaluation data
    rather than of any method, and is pooled over all checkpoints of the
    environment.}
    \label{fig:degeneracy}
\end{figure}
% [draft, compressed] DONE. The second paragraph is what remains of the
% approved paper's control section, which was a full section and is now a
% paragraph plus appendix D. It keeps the algebraic claim intact and reports
% what it buys.
A predictor that ignores its input cannot become more accurate as the horizon grows, unless the states it is asked to predict are themselves becoming more predictable. On pendulum they are. Figure~\ref{fig:degeneracy} plots the error of the dataset-mean predictor against horizon. It peaks at $6.41$ at $T{=}5$ and falls to $0.881$ by $T{=}50$, a factor of seven. Under a random policy the pendulum settles toward hanging, so the long-horizon state becomes predictable in advance of any model, and the benchmark stops separating methods at exactly the horizons where the compounding argument is most interesting. Cartpole does not degenerate, rising from $0.160$ to $2.22$ across the same range. The curve costs one line to plot and says whether a long-horizon benchmark still discriminates.

The control claim does not transfer either. The Cayley parameterization rules out divergent modes by construction, which removes one documented obstacle to a finite-cost discrete-time LQR, and that algebra holds. What it buys empirically is nothing. Solving the Riccati equation with a method that reports failure, rather than iterating the recursion and reading a gain off whatever it reaches, the unconstrained operator admits no stabilizing solution on any of the ten seeds in either environment. The Cayley operator admits one on all ten, and the resulting closed loop has spectral radius $0.9999$, a contraction so slight that a disturbance takes roughly $8{,}600$ steps to halve, against rollouts of $50$. With every mode on the unit circle, stabilizability needs all $D{=}256$ modes to be controllable through a scalar input, and the conditioning in Appendix~\ref{app:lqr} shows how far that is from holding.

% =============================================================================
